\documentclass[11pt,a4paper]{article}
\usepackage[margin=21mm,headheight=15pt]{geometry}
\usepackage[T1]{fontenc}
\usepackage{lmodern,amsmath,amssymb,booktabs,tabularx,enumitem,xcolor,fancyhdr}
\usepackage[hidelinks]{hyperref}
\definecolor{accent}{RGB}{0,114,178}
\setlength{\parindent}{0pt}
\setlength{\parskip}{6pt}
\setlist{nosep,leftmargin=*}
\pagestyle{fancy}
\fancyhf{}
\fancyhead[L]{\small Econometrics 25117 | Instructor guide}
\fancyhead[R]{\small 2026--2027}
\fancyfoot[L]{\small Private teaching notes}
\fancyfoot[R]{\thepage}
\newcommand{\heading}[1]{\section*{\color{accent}#1}}
\newcommand{\slideheading}[1]{\subsection*{\color{accent}#1}}
\newcommand{\cue}[2]{\par\noindent\textbf{#1}\quad #2\par}
\newcommand{\slideguide}[3]{\slideheading{Slide #1 | #2}\cue{Guide}{#3}}
\setlength{\emergencystretch}{1em}
\newcommand{\E}{\mathbb{E}}
\newcommand{\Var}{\operatorname{Var}}
\newcommand{\Cov}{\operatorname{Cov}}
\begin{document}
\begin{center}{\Large\bfseries Instrumental variables and 2SLS}\end{center}
\textbf{Slide route:} Lecture9v2.tex: IV, 2SLS, buttery example, cigarette application, exogeneity and relevance. Chapter 12.
\heading{Session 12 | Where does identifying variation come from?}
\textbf{Outcomes:} State relevance and exclusion/exogeneity, calculate a just-identified IV estimate, explain 2SLS, and recognize weak-instrument concerns.
\heading{100-minute route}
\begin{tabularx}{\textwidth}{@{}rX@{}}\toprule
0--10 & Retrieve simultaneity and OVB.\\
10--30 & Instrument logic and the ratio estimator.\\
30--45 & 2SLS and the lecture's buttery example.\\
45--55 & Pause and student exclusion critiques.\\
55--75 & General IV model and cigarette application.\\
75--90 & Relevance, weak instruments, exogeneity checks.\\
90--100 & Exit ticket and panel-data transition.\\\bottomrule
\end{tabularx}
\heading{Slide-by-slide script}
\slideguide{1}{Title}{Open with the question: where can credible identifying variation come from when treatment is endogenous?}
\slideguide{2}{Roadmap}{Preview IV logic, 2SLS, the buttery example, the cigarette application, relevance, and exogeneity.}
\slideguide{3}{Last lesson}{Retrieve simultaneity and ask why a strong association can still be a biased causal estimate.}
\slideguide{4}{Instrumental variables}{Define an instrument as a source of treatment variation that is relevant and excluded from the outcome equation.}
\slideguide{5}{Instrumental variables}{Use the randomized encouragement table: reduced form $=66-60=6$, first stage $=.50-.20=.30$.}
\slideguide{6}{Instrumental variables}{Calculate the Wald ratio $6/.30=20$ and explain its units.}
\slideguide{7}{Two stages least squares}{Describe first-stage prediction and second-stage outcome regression without pretending the manually generated second-stage SE is valid.}
\slideguide{8}{2SLS}{Ask what variation remains in the fitted treatment: variation induced by the instrument and controls.}
\slideguide{9}{Buttery example}{Identify the instrument, treatment, outcome, and proposed exclusion restriction before reading the numbers.}
\slideguide{10}{Buttery example}{Work the reduced-form and first-stage contrasts on the board.}
\slideguide{11}{Buttery example}{Ask students whether the instrument can affect the outcome through taste, income, or another direct channel.}
\slideguide{12}{Buttery example}{Connect the ratio calculation to a treatment effect for the variation moved by the instrument.}
\slideguide{13}{Buttery example}{State that heterogeneous effects require care: IV need not equal everyone's effect.}
\slideguide{14}{2SLS model}{Write $Y=\beta_0+\beta_1D+X'\gamma+u$ and identify the endogenous regressor.}
\slideguide{15}{2SLS model}{Write the first stage $D=\pi_0+\pi_1Z+X'\pi+v$ and explain relevance as $\pi_1\neq0$ in context.}
\slideguide{16}{Application setup}{Before using the cigarette example, name the economic question and the source of exogenous price variation.}
\slideguide{17}{Application}{Interpret the IV coefficient in the units of quantity and price, not as a generic correlation.}
\slideguide{18}{Application}{Ask whether the proposed instrument shifts demand directly. This is an exclusion question.}
\slideguide{19}{Application}{Compare OLS and IV as answers to different identifying assumptions, not as a contest where one must always be larger.}
\slideguide{20}{General IV model}{List the maintained assumptions: random sampling, relevance, exclusion/exogeneity, and sufficient variation.}
\slideguide{21}{IV exogeneity}{Draw the forbidden direct path from $Z$ to $Y$ and ask students to explain why it invalidates IV.}
\slideguide{22}{IV exogeneity}{Emphasize that a strong first stage cannot prove exclusion.}
\slideguide{23}{IV relevance}{Explain weak instruments as weak predictive power for the endogenous regressor after controls.}
\slideguide{24}{IV relevance}{Ask why weak instruments make estimates noisy and can distort usual large-sample inference.}
\slideguide{25}{IV relevance}{Interpret first-stage diagnostics as context-dependent evidence, not a magic cutoff.}
\slideguide{26}{Summary}{Have students defend or reject an instrument using separate relevance and exclusion arguments.}
\slideguide{27}{Material}{Assign one ratio calculation and one exclusion critique before moving to pooled and panel data.}

\heading{Worked example | A randomized encouragement}
\cue{Use with slides}{Lecture9v2 slides 4--7: Instrumental Variables and Two Stages Least Squares. Use the encouragement table after slide 5 to make the Wald ratio concrete, then connect it to the two-stage explanation on slides 6--7.}
Synthetic group means by encouragement $Z$:
\begin{center}\begin{tabular}{lrr}\toprule
 & $Z=0$ & $Z=1$\\\midrule
Treatment take-up $D$ & .20 & .50\\
Outcome $Y$ & 60 & 66\\\bottomrule
\end{tabular}\end{center}
\[
\hat\beta_{\rm IV}=
\frac{\bar Y_{Z=1}-\bar Y_{Z=0}}{\bar D_{Z=1}-\bar D_{Z=0}}
=\frac6{.30}=20.
\]
\cue{Explain the units}{The reduced form is six outcome units per encouragement contrast; the first stage is $.30$ treatment units per encouragement contrast. The ratio gives 20 outcome units per treatment unit.}
\cue{Assumptions}{Encouragement must shift take-up and must not affect outcomes through another channel; random encouragement supports independence. With heterogeneous treatment effects, a local average effect for compliers additionally needs monotonicity and the other usual treatment-design assumptions. Do not call it everyone's treatment effect automatically.}

\newpage
\heading{Two stages and the identifying argument}
\cue{Use with slides}{Lecture9v2 slides 14--20 for the general IV model, single-endogenous-regressor setup, application framing, and IV assumptions. Return to slides 21--26 when discussing exogeneity, relevance, and weak instruments.}
For one endogenous $X$ and exogenous controls $W$:
\[
X_i=\pi_0+\pi_1 Z_i+\pi_W'W_i+v_i,
\quad
Y_i=\beta_0+\beta_1\hat X_i+\beta_W'W_i+\text{residual}.
\]
All included exogenous controls belong in the instrument set/first stage. Multiple instruments can predict the endogenous regressor jointly.
\cue{Software caution}{The two stages explain the point estimate. The ordinary second-stage OLS standard errors from manually generated fitted regressors are not valid IV standard errors. Use an IV estimator's covariance calculation.}
\cue{Illustrative command}{With appropriately prepared variables:}
\begin{verbatim}
ivregress 2sls outcome controls (treatment = instrument), vce(robust)
\end{verbatim}
The names are placeholders; this is not a command already tested on a supplied dataset.
\heading{Weak instruments and exclusion}
If the first-stage difference were $.02$ with the same reduced form of 6, the ratio would be 300. Small denominator variation makes the estimate unstable and conventional approximations unreliable.
\cue{First-stage diagnostics}{Examine excluded-instrument strength conditional on controls. A rule of thumb is not a universal guarantee; design, estimator, and inference assumptions matter.}
\cue{Exogeneity}{A strong first stage does not prove exclusion. A transport subsidy that changes both attendance and household disposable income could affect outcomes through multiple channels. Ask students to draw the extra path.}
\cue{Overidentification}{With extra instruments, tests can detect certain violations under maintained assumptions. Failure to reject does not certify that all instruments are valid. An exactly identified model has no overidentifying restrictions to test.}
\cue{Exit answers}{Relevance concerns association with the endogenous regressor conditional on controls. Exclusion/exogeneity concern the outcome error and direct channels. Strong does not mean valid; robust does not mean weak-IV-proof.}
\cue{If short / preparation}{Work the ratio and one empirical story carefully; assign repeated application screenshots. Read the pooled/panel parts of Chapter 10 and the syllabus's Chapter 13 material as relevant.}
\textbf{Source:} Lecture9v2.tex; encouragement numbers are synthetic.

\end{document}
